% 小行星卫星（综述）
% license CCBYSA3
% type Wiki

本文根据 CC-BY-SA 协议转载翻译自维基百科\href{https://en.wikipedia.org/wiki/Minor-planet_moon}{相关文章}。

\begin{figure}[ht]
\centering
\includegraphics[width=6cm]{./figures/fe309a0bf64922bd.png}
\caption{} \label{fig_Minor_2}
\end{figure}
一个小行星卫星是指作为天然卫星环绕一颗小行星运行的天文对象。截至 2022 年 1 月，已有 457 颗已知或被怀疑拥有卫星的小行星。之所以说对小行星卫星（以及一般意义上的双天体系统）的发现十分重要，是因为通过确定它们的轨道，可以估计主星的质量与密度，从而获得通常难以通过其他方式获取的物理性质信息。

在这些卫星中，有若干在相对尺寸上非常巨大：90 号 Antiope、Mors–Somnus 与 Sila–Nunam 的尺寸约为其主星的 95\%；Patroclus–Menoetius、Altjira 与 Lempo–Hiisi 的尺寸约为 90\%，其中 Lempo–Paha 约为 50\%。已知绝对尺寸最大的小行星卫星是冥王星的卫星 Charon，其直径约为冥王星的一半。

此外，在若干遥远的天体周围，人们还发现了多个环系统（参见：Chariklo 与 Chiron 的环）。
\subsubsection{术语}
除了使用“卫星”和“月”这两个术语之外，“双星”（binary minor planet）一词有时也用于描述拥有一颗卫星的小行星，而“三级系统”（triple）用来描述拥有两颗卫星的小行星。如果其中一个天体明显更大，则称其为主星（primary），另一颗称为伴星（secondary）。

当小行星与其卫星尺寸相近时，有时使用“双小行星”（double asteroid）一词；而“binary”一词通常与两者的相对大小无关。当双小行星系统中的两颗天体大小相似时，小行星中心（MPC）将它们称作“binary companions”，而不会简单地将较小者称为卫星。

一个真正意义上的双星系统实例是 90 号 Antiope 系统，其性质于 2000 年 8 月被确认。非常小的卫星则通常称为 moonlets。

\subsection{发现历程}
在哈勃太空望远镜时代和深空探测器抵达太阳系外部之前，寻找小行星卫星的尝试主要依赖于地基光学观测。例如，1978 年，人们通过恒星掩星观测声称发现 532 号 Herculina 拥有一颗卫星。然而，之后哈勃太空望远镜的高分辨率成像并没有发现任何卫星，目前的共识认为 Herculina 并没有显著卫星。

在随后的几年中，还有其他关于小行星伴星的类似报告。当时，天文学家 Thomas Hamilton 在《Sky \& Telescope》杂志中写信指出，地球上某些看似同时形成的撞击坑（例如魁北克的 Clearwater 湖）可能是由成对、相互引力束缚的小天体撞击形成的。

同样在 1978 年，人们发现了冥王星最大的卫星 Charon；但当时冥王星仍被视为一颗主要行星。

1993 年，伽利略号探测器发现一颗小卫星 Dactyl 正环绕小行星 243 号 Ida 运行，这是首个被确认的小行星卫星。第二个例子于 1998 年在小行星 45 号 Eugenia 周围被发现。2001 年，小行星 617 Patroclus 及其等尺寸伴星 Menoetius 成为木星特洛伊群中首个已知的双小行星系统。继冥王星–Charon 系统之后，第一个被清晰分辨的海王星外双星——1998 WW31——于 2002 年被确认。

2021 年，130 号 Elektra 被发现拥有三颗卫星，使其成为目前唯一已知的四重小行星系统。
\subsubsection{多重系统}
在 2005 年，小行星 87 号 Sylvia 被发现拥有两颗卫星，使其成为首个已知的三级系统（也称三级小行星或三级小行星系统）。随后，人们又在小行星 45 号 Eugenia 周围发现了第二颗卫星。同样在 2005 年，矮行星 Haumea 被发现拥有两颗卫星，使其成为继冥王星之后第二个已知拥有多颗卫星的海王星外天体。此外，小行星 216 号 Kleopatra 与 93 号 Minerva 分别在 2008 年和 2009 年被确认为三级小行星系统。目前已知唯一的四重小行星系统是 130 号 Elektra。自最初的多重小行星系统被发现以来，后续仍不断有新的多重系统被确认。截至 2025 年，小行星中的已知多重系统总数达到 18 个（包括冥王星与 Haumea 系统）。

下表列出了所有多重系统的卫星，从冥王星开始。冥王星在其第一颗卫星于 1978 年被发现时尚未获得小行星编号。当前已知的最高重数系统为冥王星（六重系统）与 130 号 Elektra（四重系统）。
\begin{figure}[ht]
\centering
\includegraphics[width=14.25cm]{./figures/6619e07290fe4aba.png}
\caption{} \label{fig_Minor_1}
\end{figure}